% One-time imported draft. Editable native source; no Markdown regeneration.
\section{Use invariants to explain traversal independence and
  termination}\label{008-life-generations-p10-use-invariants-to-explain-traversal-independence-and-termination}

At the start of generation g\textquotesingle s sweep, current is a
complete binary board and receives no stores. Before each destination
cell, the previously visited destination cells equal the next-state rule
applied to current. Unvisited destination contents do not matter because
the counting loads never use them.

Each cell result is therefore independent of the order in which other
destination cells were visited. After the sixty-fourth result, the
entire destination is g+1. Exchanging roles preserves the
complete-current invariant. Rendering maps that state into the frame,
and publication exposes a coherent observation boundary.

The nested loops visit eight columns in each of eight rows. The outer
generation count advances once after a complete sweep and stops at the
requested zero-through-four limit. These finite counters give a
termination argument. The simulator\textquotesingle s instruction
ceiling catches broken implementations; it is not the reason the correct
loops terminate.
